---
title: "Application of the CLT: Stirlings Formula"
author: "Eugen Nesbakken"
date: "2026-03-01"
description: "A proof of Stirlings Formula for the factorial function using the Central Limit Theorem"
category: mathematics
tags: [probability]
---

\section{Introduction}

The Central Limit Theorem is one of the most important, and widely renowned results in probability theory. The statement
itself is about convergence in distribution to a Gaussian distribution.

\begin{theorem}[Central Limit Theorem]
  Let $\{X_k : k \ge 1\}$ be a sequence of i.i.d random variables with expectation $\mu < \infty$ and variance $\sigma^2 < \infty$.
  Set $S_n = \sum_{k=1}^{n} X_k$. Then
  $$
  \frac{S_n - n\mu}{\sigma\sqrt{n}} \xrightarrow{d} N(0,1)
  $$
\end{theorem}

A standard proof uses characteristic functions: the characteristic function of the standardized sum converges pointwise to $e^{-\xi^2/2}$, which characterises the standard normal.

\section{Stirlings Formula}

Stirling's formula is an asymptotic approximation of the factorial function.

\begin{theorem}[Stirling's Formula]
  $$n! \sim \left(\frac{n}{e}\right)^n \sqrt{2\pi n}$$
\end{theorem}

Where $f(n) \sim g(n)$ means $f(n)/g(n) \to 1$ as $n \to \infty$.

\section{Proof of Stirling's Formula}

\subsection{Setup}

Let $\{X_k : k \ge 1\}$ be i.i.d $\mathrm{Exp}(1)$ random variables, so $E[X_k] = 1$ and $\mathrm{Var}(X_k) = 1$. Set
$$
S_n = \sum_{k=1}^n X_k \sim \Gamma(n, 1)
$$
so that $S_n$ has density
$$
f_{S_n}(x) = \frac{x^{n-1} e^{-x}}{(n-1)!}\,\mathbb{1}_{\{x \ge 0\}}.
$$
Define the standardized sum $Z_n = (S_n - n)/\sqrt{n}$, which has density
$$
g_n(z) = \sqrt{n}\, f_{S_n}(n + z\sqrt{n}).
$$
Evaluating at $z = 0$:
$$
g_n(0) = \sqrt{n}\cdot\frac{n^{n-1}e^{-n}}{(n-1)!}.
$$
Multiplying both sides by $n$ shows that Stirling's formula is equivalent to
$$
g_n(0) \;\to\; \frac{1}{\sqrt{2\pi}}.
$$
We prove this limit rigorously using characteristic functions.

\subsection{Characteristic Function of $Z_n$}

The characteristic function of $X_k \sim \mathrm{Exp}(1)$ is
$$
\varphi_X(\xi) = E[e^{i\xi X}] = \frac{1}{1 - i\xi}, \qquad \xi \in \mathbb{R}.
$$
By independence,
$$
\varphi_{Z_n}(\xi)
= e^{-i\xi\sqrt{n}}\,\varphi_X\!\left(\frac{\xi}{\sqrt{n}}\right)^n
= e^{-i\xi\sqrt{n}}\left(1 - \frac{i\xi}{\sqrt{n}}\right)^{-n}.
$$
Expanding the logarithm for fixed $\xi$ as $n \to \infty$:
$$
-n\ln\!\left(1 - \frac{i\xi}{\sqrt{n}}\right)
= -n\left(-\frac{i\xi}{\sqrt{n}} - \frac{\xi^2}{2n} + O(n^{-3/2})\right)
= i\xi\sqrt{n} - \frac{\xi^2}{2} + O(n^{-1/2}).
$$
Therefore
$$
\varphi_{Z_n}(\xi) = \exp\!\left(i\xi\sqrt{n} - i\xi\sqrt{n} - \frac{\xi^2}{2} + O(n^{-1/2})\right) \;\longrightarrow\; e^{-\xi^2/2}
$$
pointwise for each $\xi \in \mathbb{R}$.

\subsection{Dominated Convergence and Fourier Inversion}

Since $g_n$ is a continuous density, the Fourier inversion formula gives
$$
g_n(0) = \frac{1}{2\pi}\int_{-\infty}^{\infty} \varphi_{Z_n}(\xi)\,d\xi.
$$
We justify passing the limit $n \to \infty$ inside the integral via dominated convergence. The modulus of $\varphi_{Z_n}$ is
$$
|\varphi_{Z_n}(\xi)| = \left(1 + \frac{\xi^2}{n}\right)^{-n/2}.
$$
The function $n \mapsto (1 + \xi^2/n)^{n/2}$ is increasing in $n$ (since its logarithm $(n/2)\ln(1+\xi^2/n)$ has non-negative derivative in $n$), so for all $n \ge 4$:
$$
|\varphi_{Z_n}(\xi)| \;\le\; \left(1 + \frac{\xi^2}{4}\right)^{-2}
$$
which is integrable over $\mathbb{R}$. The Dominated Convergence Theorem therefore gives
$$
g_n(0) = \frac{1}{2\pi}\int_{-\infty}^{\infty} \varphi_{Z_n}(\xi)\,d\xi
\;\longrightarrow\;
\frac{1}{2\pi}\int_{-\infty}^{\infty} e^{-\xi^2/2}\,d\xi
= \frac{\sqrt{2\pi}}{2\pi} = \frac{1}{\sqrt{2\pi}}.
$$

\subsection{Conclusion}

Recalling that $g_n(0) = \sqrt{n} \cdot n^{n-1}e^{-n}/(n-1)!$, we have shown
$$
\sqrt{n}\cdot\frac{n^{n-1} e^{-n}}{(n-1)!} \;\longrightarrow\; \frac{1}{\sqrt{2\pi}}.
$$
Multiplying through by $n$:
$$
n! \;\sim\; \sqrt{2\pi n}\cdot n^n \cdot e^{-n} = \sqrt{2\pi n}\left(\frac{n}{e}\right)^n. \qquad \blacksquare
$$

\begin{remark}
 The exact form of the characteristic function of Exp(1) and its pointwise convergence to $e^{-\xi^2/2}$ is essential here. We cannot deduce based on the densities, but the CDFs since CLT does not imply pointwise convergence of the densities. The $\sqrt{2\pi}$ in Stirling's formula is precisely the $L^1$ norm of $e^{-\xi^2/2}$, i.e.\ the normalising constant of the Fourier transform of the Gaussian.
\end{remark}
